\subsubsection*{The residual coefficient of the signature determinant}
\label{subsec:residual-twist}

The signature lattice now gives a square matrix. To control its determinant at a specialization, we need a coefficient detectable from the mod-\(\ell\) Selmer classes and the Frobenius elements chosen for them.

Keep the notation of Section~\ref{subsec:signature-selmer} and write
\(W=T\otimes k\), \(W'=W^\vee(1)\). Conjugation of a group
element by complex conjugation is denoted by \(h\mapsto h^c\).
The good-signature duality gives an integral isomorphism of
\((T^\vee(1))^c\) with \(T\otimes\epsilon^t\), where \(\epsilon\)
is cyclotomic. The tame marking gives an inertia element \(z\)
acting on \(W\) by \(d\), with \(d^2\ne1\), and on \(W'\) by
\(d^{-1}\). At \(z^c\) the two scalars are reversed.
Both modules are absolutely irreducible. Let \(S,S'\) be the
residual Selmer groups for these modules at empty auxiliary
level, with finite local conditions at \(\ell\) defined by
the integral lattices.

\begin{lemma}[Residual Selmer dimensions]
For the residual primal and dual Selmer spaces of a good signature over its residue field \(k\),
\[
 \dim_k S=\dim_k S'.
\]
\end{lemma}
\begin{proof}
At \(v\mid\ell\), the cup of two integral finite cocycles has
zero rational invariant by Bloch--Kato orthogonality. The local
invariant map embeds integral \(H^2\) in rational \(H^2\) on the
scalar Tate twist, hence the reductions are orthogonal. For the
lattice of \(W\), Section~\ref{subsec:signature-selmer} gives
\[
 \dim L_v(k)-\dim H^0(E_v,W)
   =\dim H^1_f(E_v,V)-\dim H^0(E_v,V).
\]
The right side is the dimension of the de Rham quotient by
\(\operatorname{Fil}^0\). Add this equality to its analogue for
the dual lattice. At each embedding the two quotient dimensions
sum to \(r\) by the dual Tate filtration. Local duality identifies the invariants of the Tate dual with the
dual of \(H^2(E_v,W)\). Thus the local Euler characteristic reads
\[
 \dim L_v(k)+\dim L'_v(k)
 = [E_v:\mathbb Q_\ell]r+\dim H^0(E_v,W)+\dim H^2(E_v,W)
 =\dim H^1(E_v,W).
\]
The orthogonal subspaces therefore have complementary dimensions.

At the type and distinguished places
away from \(\ell\) the inertia test of the signature excludes
invariants on the residual modules: an
inertia element acts with no eigenvalue one. The dual has
the inverse eigenvalues there. The local Euler characteristic
and duality give \(H^1=0\). Away from these places and
\(\ell\), the unramified subgroups are exact annihilators.
Global duality (the finite-module Selmer order formula) now
gives \(|S|/|S'|=1\): the global invariant factors are one,
the complex factors contribute \(|W|^{-f}\), and the
\(\ell\)-adic factors contribute \(|W|^f\) by
\eqref{eq:signature-local-size}. \cite{bloch-kato,milne-duality}

\end{proof}
\paragraph{Forms specified by Frobenius elements.}
There is an open normal subgroup \(J\subset H\), stable under
\(c\), on which the residual representations and the action
on \(\mu_{2L}\) are trivial and which contains all commutators
\([g,z]\). Intersect \(H\) with the cyclotomic and residual
kernels and with the conjugate of the intersection. The
commutators lie in these groups because \(H\) contains
inertia at the distinguished pair, cyclotomic is abelian,
and the marked elements act by scalars. The same construction
works simultaneously for finitely many signatures. If needed
one may replace each residual kernel here by the kernel modulo
\((\ell^{\nu_\ell})\mathfrak m_{O_0}\): the distinguished
inertia acts by scalars integrally on \(T\).

For \(h\in J\) put
\[
       B_h(x,y)=x(h)\cdot y(h)-x(h^c)\cdot y(h^c)
            \qquad (x\in S,\ y\in S').
\]
\begin{lemma}[Corank of the Frobenius pairing]
\label{lem:residual-frobenius-forms}
For every good signature there is a linear combination \(B\) of the second-half Frobenius forms for which
\[
 \dim\ker\bigl(S\longrightarrow(S')^*,\ x\longmapsto B(x,\cdot)\bigr)\le \mathbf1_{\nu_\ell>0}.
\] The same second-half primes can be
used for the finitely many signatures. Before combining their
lists, each requires at most \((\dim S+\dim S')^2\) elements of \(J\).
\end{lemma}
\begin{proof}
The values are independent of the cocycle representatives.
Restriction to \(J\) is injective on cohomology. In fact a
cocycle \(x\) vanishing there factors through a quotient in
which \(z\) is central. The two expressions for \(x(gz)=x(zg)\)
give
\[
 (d-1)x(g)=(g-1)x(z),\qquad
 x(g)=(g-1)\bigl((d-1)^{-1}x(z)\bigr).
\]
Thus \(x\) is a coboundary. On the dual use \(d^{-1}\) in the
same identity; both denominators are nonzero. For independent classes
\(x_1,\ldots,x_a\) with values in \(W\), the vectors
\((x_j(h))_j\) span \(W^a\) over \(k\) as \(h\) varies in
\(J\). Let \(U\) be the span of the evaluation tuples in \(W^a\). It is stable by conjugation. If \(U\) were proper, semisimplicity and absolute irreducibility would give scalars, not all zero, with
\[
 \sum_{j=1}^a c_j x_j(h)=0\qquad(h\in J).
\]
Indeed any map \(W^a\to W\) on this isotypic space is a linear combination of coordinate projections. Restriction injectivity rules it out.

If \(W\not\simeq (W')^c\), then
\(W,W',W^c,(W')^c\) are pairwise inequivalent: the scalars
at \(z\) separate the two groups and conjugation carries the
remaining inequivalence to its partner. For bases of \(S,S'\)
the simultaneous evaluations of these four lists on \(h\)
span the direct sum of their four full isotypic spaces. This
follows from the preceding projection test on each simple
type. To show that the \(B_h\) span all bilinear forms, suppose
\[
 \sum_{i,j}a_{ij}B_h(x_i,y_j)=0
\]
for every \(h\in J\).

Write the four evaluation blocks as \(v_i,w_j,\tilde v_i,\tilde w_j\),
so the relation says that every evaluation tuple is a zero of
\[
 Q(v,w,\tilde v,\tilde w)=\sum_{i,j}a_{ij}
      (v_i\cdot w_j-\tilde v_i\cdot\tilde w_j).
\]
Evaluation on \(J\) is additive, so the bilinear form
\(B_Q(s,t)=Q(s+t)-Q(s)-Q(t)\) vanishes on all pairs of evaluation
tuples and hence on their full linear span. Fix \(i,j\); give
\(s\) just the entry \(v\) in the \(i\)-th first block and
\(t\) just the entry \(w\) in the \(j\)-th second block. Then
\begin{equation}\label{eq:frobenius-polarization-block-test}
 0=B_Q(s,t)=a_{ij}\,v\cdot w .
\end{equation}
Choosing a pair with nonzero pairing in \eqref{eq:frobenius-polarization-block-test} forces \(a_{ij}=0\), including
in characteristic two.

If \(W\simeq(W')^c\), choose an invertible intertwining
matrix \(A\), so that
\(W(g^c)^t A W(g)=\bar\epsilon(g) A\), with superscript
\(t\) on a matrix denoting transpose. Transposing the same
identity at \(g^c\) shows \(A^t=\lambda A\) by Schur's lemma.
Here \(\lambda^2=1\) and \(\lambda^r=1\) by determinants, so
\(A^t=A\) since \(r\) is odd. Transport \(S'\) to a subspace
\(S_2\) of \(H^1(E,W)\) by sending \(y\) to \(A^{-1}y^c\).
Choose a basis of \(S+S_2\) and evaluate it at \(h,h^c\).
The two evaluations span the two full isotypic blocks together
because \(z\) separates \(W\) and \(W^c\).

In this basis the
form has matrix entries
\[
     x_i(h)^t A x_j(h^c)-x_i(h^c)^t A x_j(h).
\]
A relation with coefficients \(a_{ij}\) depends on
\(a_{ij}-a_{ji}\). Polarization and testing on vectors in
the two separate blocks forces \(a_{ij}=a_{ji}\). Conversely
these symmetric relations annihilate the displayed alternating
matrices, including their diagonals. Thus the forms span the restrictions of alternating forms on
\(S+S_2\). Set \(J_0=S\cap S_2\) and take complements
\(S=J_0\oplus U\), \(S_2=J_0\oplus U'\). These complements
have the same dimension. With dual bases for the pairing
of \(U\) with \(U'\), choose the restriction \(S\times S_2\)
with block matrix
\[
 \begin{pmatrix}J_{\dim J_0}&0\\0&I_{\dim U}\end{pmatrix},
 \qquad
 J_d=\bigoplus_{i=1}^{\lfloor d/2\rfloor}
       \begin{pmatrix}0&1\\-1&0\end{pmatrix}
       \quad\text{with a final zero block if \(d\) is odd}.
\]
The opposite off-diagonal pairing on \(U'\times U\) is the
negative transpose, so this matrix extends to an alternating
form on \(S+S_2\). Its left corank is \(d\bmod 2\), including
in characteristic two.

Consequently there is a linear combination of the forms
with left corank at most \(e\), where \(e=0\) if
\(\nu_\ell=0\) and \(e=1\) otherwise. Indeed an equivalence
\(W\simeq(W')^c\) would imply
\[
 \bar\epsilon^{tr}=1
\]
by taking determinants in the signature duality, which
implies \(\nu_\ell>0\) by \eqref{eq:signature-cyclotomic-minimum}.
The number of elements \(h\) needed is at most
\((\dim S+\dim S')^2\), independently of \(J\) and \(M\).
Choose such elements for each good signature and concatenate
the lists, padding with identities if necessary.

For the finite list \(\mathcal Z\) of restrictions of the basis cocycles and their conjugates, use the common kernel
\[
 J_0=\bigcap_{z\in\mathcal Z}\ker(z|_J).
\]
The conjugation formula gives an open normal subgroup of \(J\).
Apply Section~\ref{subsec:aux-coordinates} prescribing these
elements for the second half. On replacing \(h\) by a
Frobenius conjugate to it modulo this subgroup the matrix
\(B_h\) changes only by the nonzero cyclotomic scalar of the
conjugating element. Absorb it into the coefficients. Thus at
the second-half Frobenius elements \(\phi_i\) we have, for each signature,
coefficients \(b_i\) such that \(\sum_i b_i B_{\phi_i}\) has
the stated corank. The same primes work for all signatures.

\end{proof}
\begin{lemma}[A residual coefficient from the Selmer pencil]\label{lem:residual-selmer-coefficient}
Use the two halves of auxiliary places and the Frobenius forms
of Lemma~\ref{lem:residual-frobenius-forms}, with \(R\) places
in each half and \(L>4rR+2\). When \(\nu_\ell=0\), every polynomial determinant lift \(D\) satisfies
\[
 \exists\alpha:\quad |\alpha|\le4rR,\qquad [X^\alpha]D\in O^\times.
\] For positive \(\nu_\ell\) the corresponding
unit coefficient is obtained at this stage in an adjugate
entry, for a nonempty square presentation. An empty presentation
already has determinant one.
\end{lemma}

\begin{proof}
\proofstep{Use the inertia and Frobenius evaluations in the relaxed residual space.}
Let \(R\) denote the common length of each half. At the first
half the untwisted residual module is unramified with Frobenius
scalar \(d\) or \(d^{-1}\); both differ from one. The
cyclotomic character is trivial there modulo \(\ell\).
Thus \(H^0,H^2\), and by the local Euler characteristic also
\(H^1\), vanish. At each second-half prime the Frobenius is
in \(J\), as is inertia. Fix a tame \(\ell\)-generator
\(\gamma_i\) with \(d_{P'_i}(\gamma_i)=1\). This is possible
since the character on inertia is onto. In the following
evaluation pairs, use at the conjugate place the transports
\(\gamma_i^c,\phi_i^c\). Superscript \(c\) on group elements
continues to denote conjugation.

Write \(S_Q\) for the relaxed residual Selmer space and
\[
 D=(k^r\oplus k^r)^R.
\]
Define \(I:S_Q\to D\) by evaluation
at \((\gamma_i,\gamma_i^c)_i\). Its kernel is \(S\), since
on the trivial residual module a homomorphism on inertia
factors through its pro-\(\ell\) tame quotient. The first-half
places contribute nothing. For \(x\in S_Q\), \(y\in S'\),
reciprocity gives units \(a_i\in k^\times\), independent of
\(x,y\), such that
\begin{equation}\label{eq:residual-reciprocity}
 \sum_i a_i\{x(\gamma_i)\cdot y(\phi_i)
             -x(\gamma_i^c)\cdot y(\phi_i^c)\}=0 .
\end{equation}
Indeed the cup invariant at an unramified pair of classes is
zero, and at the primes over \(\ell\) it vanishes by the local
orthogonality above.

At a second-half place the second class
is unramified. The invariant on a tame character of
\(\gamma_i\) paired with an unramified character of \(\phi_i\)
is a nonzero scalar: the two unramified classes cup to zero
and local Tate duality is nondegenerate on the two-dimensional
space for the trivial residual line. Bilinearity gives the
formula for vectors. The transport to the conjugate place
reverses the chosen identification of \(\mu_\ell\) with the
prime field, since \(c\) acts by \(-1\) there. This explains
the common scalar and the minus sign. The sum of local
invariants vanishes by global reciprocity.

Define
\[
 F:S_Q\to D
\]
by taking the pairs of Frobenius
values of each class multiplied by \(b_i/a_i\). The two necessary linear conditions on a lift are
\[
 Ix=0,\qquad Fx=Ix_1\quad(x_1\in S_Q).
\]
The first gives \(x\in S\) and \eqref{eq:residual-reciprocity} applied
to \(x_1\) puts \(x\) in the left radical of
\(\sum b_i B_{\phi_i}\). Consequently
\begin{equation}\label{eq:residual-pencil}
  \dim\{x\in\ker I: Fx\in\operatorname{im}I\}\le e,
  \qquad \dim S_Q\le 2rR+e .
\end{equation}
For the second inequality apply rank-nullity to
\(\ker I\to D/\operatorname{im}I\) induced by \(F\).

We give the specialization producing a coefficient of bounded
degree; setting all group variables to the same variable would
lose the coefficients \(b_i/a_i\). Put \(K=k((s))\).
In the polynomial presentation matrix reduce modulo the
uniformizer and set \(X_j=\alpha_j Z\), \(Y_i=\beta_i Z\),
where
\[
  \beta_i=s\, b_i/a_i,\qquad
  \alpha_j=1-\sum_i\beta_i\,\overline{d_{P'_i}(\phi_j)} .
\]
Write \(P(Z)\) for this square matrix over \(K[Z]\).
Modulo \(Z^2\) the substitution factors through the group ring,
since \(L\) is a positive power of \(\ell\). On the second-half
decomposition group the coefficient of \(Z\) in the inverse
universal character is \(-1\) at \(\phi_i\) and \(-\beta_i\)
at \(\gamma_i\), by \eqref{eq:aux-coordinate}. At the
conjugate place both signs reverse.

A first-order lift of \(v\in\ker P(0)\) is a vector \(v+Zv_1\) satisfying
\[
 P(Z)(v+Zv_1)\equiv0\pmod{Z^2}.
\]
Naturality of the square presentation gives
a Selmer cocycle in the dual numbers lifting its residual
class \(x\in K\otimes_k S_Q\). This identification commutes
with taking its constant term. Indeed additive functoriality
gives the assertion for finite sums of \(k\), and each class
over a discrete extension field is supported on a finite
coefficient subspace.

At each second-half place write the
cocycle as \(x+Zx_1\), and the character as \(1+Z\lambda\).
The tame relation
\[
 \phi\gamma\phi^{-1}=\gamma^q,
\]
where
\(q=\mathrm NP'_i\), gives
\[
             \lambda(\phi)x(\gamma)=\lambda(\gamma)x(\phi).
\]
For this calculation the crossed term of the first-order
cocycle at \(g,h\) is \(\lambda(g)x(h)\). In its value at
\(\gamma^q\) the coefficient \(\binom q2\) vanishes in \(k\)
because \(q\equiv1\bmod 2L\); the coefficient \(q\) is one.
The prime-to-\(\ell\) inertia kernel has zero image in the
successive additive groups used in this calculation.
Thus \(Ix=sFx\), including at the conjugate places.

The subspace of \(K\otimes S_Q\) satisfying this equation
has dimension at most \(e\). Intersect it with the lattice
\(S_Q[[s]]\). This saturated sublattice has reduction of
dimension equal to its rank and may be given a basis with
independent constant terms by elimination over \(k[[s]]\).
Coefficient comparison in \(Ix=sFx\), with \(x=x_0+sx_1+\cdots\), gives
\[
 Ix_0=0,\qquad Ix_1=Fx_0.
\]
Now apply
\eqref{eq:residual-pencil}. This argument can be used over
any finite residue field, including \(\mathbb F_2\).

\proofstep{Use the first-order kernel to bound the number of invariant factors of order at least two.}
The evaluations have bounded the kernel of the first-order matrix. Reading the same matrix through its invariant factors converts that kernel bound into the restriction on the determinant's low-degree coefficients.

Take Smith form over \(K[[Z]]\), with orders
\(0\le a_1\le\cdots\le a_d\le\infty\), where a zero
diagonal entry has order \(\infty\). Reduce the kernel of the \(i\)-th diagonal entry on \(K[[Z]]/(Z^2)\)
modulo \(Z\). Its image in \(K\) is
\[
 \begin{cases}
 0,&a_i=0\text{ or }a_i=1,\\
 K,&2\le a_i\le\infty.
 \end{cases}
\]
For \(a_i=1\) its kernel is precisely the ideal \((Z)\), whose
residue in \(K\) is zero.
The dimensions just computed give
\[
 \#\{i:a_i>0\}=h=\dim S_Q\le2rR,
 \qquad \#\{i:a_i\ge2\}\le e.
\]
If \(e=0\), all Smith exponents are zero or one. Therefore
\[
 \operatorname{ord}_Z\det P=\sum_i a_i=\#\{i:a_i=1\}=h .
\]
If \(e=1\) and \(d>0\), the ideal generated by the submaximal
minors has order
\[
 \sum_{i=1}^{d-1}a_i\le\max(0,h-1)\le2rR.
\]

At least one adjugate entry has that order. For an empty matrix
the determinant is already one. Both conclusions suffice with
the degree bound \(4rR\). For \(L>4rR+2\) the
choice of polynomial lift cannot change any reduced coefficient
in these degrees: the coordinate relations reduce to
\(X_j^L=Y_i^L=0\). A nonzero coefficient of the line
specialization requires a nonzero coefficient of the same total
degree before substitution. Hence every lift of the determinant
has a unit coefficient of bounded degree when \(\nu_\ell=0\).
For \(\nu_\ell>0\) the conclusion at this stage concerns one
adjugate entry. Its conversion to a determinant with valuation
at most \(\nu_\ell\) requires one further arithmetic step.
\end{proof}

